
In this work, we presented the first calculation of the pion gluon PDF from lattice QCD and studied its pion-mass and lattice-spacing dependence using the pseudo-PDF approach.
We employed clover valence fermions on ensembles with $N_f = 2+1+1$ highly improved staggered quarks (HISQ) at two lattice spacings ($a\approx 0.12$ and 0.15~fm) and three pion masses (220, 310 and 690~MeV).
These ensembles allowed us to probe the dependence of the pion gluon PDF on pion mass and lattice spacing.
In both cases, the dependence appears to be weak compared to the current statistical uncertainty.

We investigated the systematics associated with the functional form used in the reconstruction fits as well as the systematics caused by neglecting the quark contribution in the matching.
The effect of the assumed gluon PDF fit form was investigated by using various forms, which are all commonly used or proposed in other PDF works.
We observe large effects changing the fit to $xg(x, \mu)/\langle x \rangle_g$ from 1- to 2-parameter form but convergence at 3 parameters.
This implies the 2-parameter fits are sufficient for our calculation, and our finial pion gluon PDF results are presented using the 2-parameter fit results.
We used the pion quark PDF from xFitter to make an estimation of the quark contribution to the pion gluon RITD.
We found the systematic errors it contributed are smaller than $10\%$ of the statistical errors.

Our pion gluon PDF for the lightest pion mass is consistent with JAM for $x>0.2$ and with xFitter for $x>0.5$ within uncertainty, as shown in our final comparison plots of the pion gluon PDF.
We also studied the asymptotic behavior of the pion gluon PDF in the large-$x$ region in terms of $(1-x)^C$.
$C>3$ is implied from our study at two lattice spacings and three pion masses.
The future study of the pion gluon PDF from the lattice QCD with improved precision and systematic control when combined in global-fit analyses with the results of
anticipated experiments~\cite{Denisov:2018unj,Arrington:2021biu,Anderle:2021wcy,Denisov:2018unj} will provide best determination of the gluon content within the pion.
